\def\ChapterCode{ABD}
\chapter
[\CO{[ABD]} Abdominale Erkrankungen]
{Abdominale Erkrankungen}
\ChapterIntro
{%
Der Bauch ist eine echte Herausforderung. Er beinhaltet
viele unterschiedliche Organe und Strukturen welche weh tun
und erkranken können. Viele dieser, oft auch mit starken
Schmerzen verbundenen Störungen sind nicht übermäßig
schlimm, einige sind jedoch sehr gefährlich oder können
rasch gefährlich werden.


}
{%
\begin{description}
  \item[Maintainer:] Sebastian Gabriel
  \item[Autoren:] Diverse
  \item[Reviewer:] Standard-Reviewprozess
  \item[Version:] {\DocVersion} ({\DocVersionInfo})
  \item[SHA1:] {\DocGitCommit}
\end{description}

{\TxTeilkapitelAASS}
}


\clearpage

\section{Allgemeines}

\TextSlide{Anamnese und Untersuchung}{%
Es ist wichtig
typische Symptome zu erheben und Alarmzeichen sofort zu
erkennen.
Wesentlich ist bei allen unklaren Bauchbeschwerden die
Erhebung einer exakten \I{Stuhlanamnese}:

\begin{itemize}
  \item Auffälligkeiten? (Geruch, Farbe, Konsistenz, Frequenz)
  \item Blutbeimengungen
  \item Wann erfolgte der letzte Stuhlgang?
  \item Winde?
  \item Genaue Beschreibung von Harn und eventuell Erbrochenem.
  \item Bei Frauen im gebärfähigem Alter Regelanamnese.
\end{itemize}

Der Patient muss bis zur endgültigen Klärung der Situation
\EE{nüchtern} belassen werden!
Wenn es der Patient toleriert, muss eine Betrachtung und
Abtastung des \so{unbekleideten} Abdomens beim liegenden
Patienten erfolgen wie unter
\ref{topic:untersuchung-abdomen-abtasten} beschrieben.
}{%
\begin{itemize}
  \item Stuhl
  \begin{itemize}
    \item Wann letzter?
    \item Auffälligkeiten
    \item Farbe, Konsistenz
  \end{itemize}
  \item Erbrochen?
  \item Blut?
  \item Regel?
\end{itemize}
}{}{}{}

\subsection{Der schmerzende Bauch}

\TextSlide{Der Bauch -- unendliche Weiten}
{%
Gleich vorweg: In vielen
Fällen wird der Grund des Bauchschmerzes nicht mit
hinreichender Sicherheit feststellbar sein und es bietet
sich keine sichere Verdachtsdiagnose an. Statt einer
Diagnose wird dann das \EE{Beschwerdebild beschrieben} (Ort
des Schmerzes, Art des Schmerzes, Dauer, Verlauf,
Austrahlung, Ergebnis der abdominellen Tastuntersuchung:
Abdomen weich oder hart,
Druckschmerz, Loslassschmerz). Dennoch muss man bei jedem Patienten versuchen, so gut als
möglich die in Frage kommenden Krankheitsbilder einzugrenzen
und gefährliche Erkrankungen nach Möglichkeit auszuschließen. Im
folgenden sind einige häufige Krankheitsbilder beschrieben,
die \E{erkennbar}  und relativ häufig sind.
}{
\begin{itemize}
  \item In vielen Fällen Krankheitsbild unklar
  \item Im Zweifel: Beschreibung statt Diagnose
\end{itemize}
}{}{}{}

\Text{Beispiele}
{für Beschreibungen:
\begin{quote}
{\sffamily\bfseries Diagnose:} {\ttfamily Unklarer
Bauchschmerz: Kolikartige Schmerzen im linken Unterbauch seit 1/2h }

{\sffamily\bfseries Anmerkungen:} {\ttfamily keine
Ausstrahlung, Druckschmerz im li. Unter\-bauch, Abdomen sonst 
weich.}
\vspace{0.5em}

{\sffamily\bfseries Diagnose:} {\ttfamily Unklarer
Bauchschmerz: Brennend-stechender Schmerz im mittleren
Oberbauch seit 2h mit Ausstrahlung in den Rücken}

{\sffamily\bfseries Anmerkungen:} {\ttfamily Abdomen weich,
keine Druck- oder Loslassschmerzen}
\vspace{0.5em}

{\sffamily\bfseries Diagnose:} {\ttfamily Unklarer
Bauchschmerz: „Dumpfer“ Schmerz im linken und rechten
Oberbauch seit gestern. }

{\sffamily\bfseries Anmerkungen:} {\ttfamily Abdomen
weich, keine Druck- oder Loslassschmerzen}
\end{quote}
}{

}{}{}{}




% M %%%%%%%%%%%%%%%%%%%%%%%%%%%%%%%%%%%%%%%%%%%%%%%%%%% M %
% Abdominalerkrankung
\ProcedurePrintDefault{CKXXXX1C}

% \clearpage % TEMP temp clearpage
\section{Häufige und gut erkennbare Erkrankungen}



\subsection{Das Akute Abdomen}
\label{topic:akutes-abdomen}
\index{Abdomen!akutes|see{Akutes Abdomen}}





\Text{\TxBeschreibung}
{\DefinitionInPara{Das \T{Akutes Abdomen} ist eine akute
Erkrankung des Bauchraumes mit plötzlichen, starken
\EE{Bauchschmerzen} und
\EE{bretthartem Bauch}, welche im weiteren Verlauf
lebensbedrohlich ist. Es sind verschiedene Ursachen möglich.}
\Commentary[Akutes Abdomen, Definition]{Die Definition des
Akuten Abdomens ist oft nicht ganz eindeutig: Einerseits kann man 
darunter jede plötzlich auftretende
Baucherkrankung
verstehen \citep{Longmore2007,Renz-Polster2006}. 
Allerdings ver\-stehen viele Leute darunter eher eine
bedrohliche Bauch\-erkrankung, bei der man akut (schnell)
handeln muss (Ab\-wehrspannung, harter Bauch). In der
Notfallmedizin halten wir uns an letztere Definition.}
Die zu Grunde liegenden Erkrankungen kann
plötzlich entstehen, oder sich über längere Zeit entwickeln
und erst im fortgeschrittenen Stadium die Symptome des
akuten Abdomens auslösen. Ein akutes Abdomen erfordert eine
rasche (spitalsmäßige) Behandlung.
}
{%
\begin{itemize}
  \item Akute Erkrankung des Bauchraumes mit
  \begin{itemize}
    \item Bauchschmerzen
    \item Brettharter Bauchdekce
  \end{itemize}
  \item Verschiedene Ursachen
\end{itemize}
}{}{}{}

\TextSlide{\TxSymptome}
{%
Der Patient liegt meistens im Bett. Oft ist eine
Schonhaltung (angezogene Beine, gekrümmte Körperhaltung) zu
beobachten. Er klagt über \EE{Schmerzen} im Bauch, diese können
als diffus (im gesamten Bauch) oder auf ein Bauchsegment
beschränkt angegeben werden.

Die \EE{harte Bauchdecke} (\enquote*{\E{bretthart}}) ist
neben den Schmerzen das Leitsymptom des Akuten Abdomens.
Es besteht oft Übelkeit oder auch Erbrechen. Der Kreislauf
kann kompensiert sein, oder Zeichen von Schwäche zeigen, bis
hin zu einer Schocksymptomatik (Blutverlust, Sepsis).
Weitere Symptome sind abhängig von der zu Grunde liegenden
Erkrankung.

Diese Auflistung der Symptome ist in der Praxis nicht
unbedingt vollständig bei einem Patienten nachzuweisen. Die
Intensität und die Anzahl der angeführten Symptome ist
einerseits vom Allgemeinzustand, vom Lebensalter, von
anderen vorbestehenden Erkrankungen des Patienten und der
Dauer des aktuellen Zustandes abhängig.

Der Patient ist nicht unbedingt unmittelbar vital
bedroht. Wesentlich nach Eintreffen beim Patienten ist die
Beantwortung der Fragen \Speech{\enquote{Wie schlecht, bzw.
wie gut geht es dem Patienten? Ist der Patient \uline{akut}
vital bedroht?}}.

} 
{
\begin{itemize}
  \item \EE{Schmerz}: diffus (ganzer Bauch) oder
  lokalisierbar
  \item \EE{Harte  Bauchdecke} (bretthart)
  \item Schonhaltung  (angezogene Beine, Patient
  krümmt sich)
  \item Übelkeit, Erbrechen
  \item Evtl. Stuhl-, Windverhalten
  \item Evtl. Fieber
  \item Evtl. Schock
\end{itemize}
}{}{}{}




\TextSlide{Ursachen}
{%
Grundsätzlich kommt fast jede Erkrankung oder Verletzung in
Betracht die Strukturen und Organe im Bauchraum betrifft,
entsprechend vielfältig sind auch die Ursachen: 
\begin{itemize}
  \item \EE{Entzündungen} im Bauchraum (Bauchfellentzündung
  (Peritonitis), Blinddarmentzündung (Appendizitis),
  Entzündung der Bauchspeicheldrüse  (Pankreatitis),\ldots)
  \item \EE{Mesenterialinfarkt}
  \item \EE{Perforation} eines Hohlorganes
  (Darmverschluss (Ileus), Blinddarmdurchbruch, \ldots)
  \item \EE{Blutungen} in den Bauchraum
  (Bauchaortenaneurysma, rupturierte Eileiterschwangerschaft, \ldots)
  \item andere abdominelle Erkrankungen
  \item Gynäkologische Erkrankungen
  \item \ldots
\end{itemize}
}{
\begin{itemize}
  \item Entzündungen (Bauchfell, Blinddarm, Pankreatitis,
  \ldots)
  \item Perforation eines Hohlorganes
  \item Blutungen (Aneurysma, Eileiterschwangerschaft,
  \ldots)
  \item Gynäkologische Erkrankungen
  \item u.\,v.\,a.\,m. \ldots
\end{itemize}
}{}{}{}




% M %%%%%%%%%%%%%%%%%%%%%%%%%%%%%%%%%%%%%%%%%%%%%%%%%%% M %
% Akutes Abdomen
\ProcedurePrintDefault{CR10000C}


\subsection{Darmverschluss}
\index{Darmverschluss}%
\label{topic:ileus}%


\Text{\TxBeschreibung}
{%
\DefinitionInPara{Als \T{Darmverschluss}
bezeichnet man die Unterbrechung der normalen Darmpassage. 
\Lang{Term.} \T{Ileus} (vollständig), \Z{\T{Subileus} (unvollständig)}}
Der Verschluss kann mechanisch bedingt sein (durch
Hindernisse wie Verwachsungen oder Tumore), oder durch eine
Lähmung der Darmmuskulatur (Medikamente, Gifte, Entzündung
\ldots) verursacht sein. Bestimmte Schmerzmittel
\Z{(Opiate)}, welche oft bei Krebspatienten eingesetzt
werden, sind bekannt dafür Darmlähmungen zu erzeugen. 
}{%
\begin{itemize}
  \item Mechanische Blockade oder Lähmung
\end{itemize}
}{}{}{}


\TextSlide{\TxSymptome}
{%
Typisch für den Darmverschluss ist die mangelhafte
Stuhlproduktion. Es kommt zu diffusen Bauchschmerzen und einem
Blähbauch, sowie zu Übelkeit und Erbrechen. In seltenen
Fällen kann es auch zum Erbrechen von zurück
gestautem Kot kommen \Z{(\I[Miserere|Z]{Miserere})}.

\bigskip

\SlideCompensationNoParskip{1}
}{
\begin{itemize}
  \item Diffuse Bauchschmerzen
  \item Stuhlverhalten
  \item Übelkeit, Erbrechen (selten: Koterbrechen)
\end{itemize}
}{}{}{}



% M %%%%%%%%%%%%%%%%%%%%%%%%%%%%%%%%%%%%%%%%%%%%%%%%%%% M %
% Darmverschluß
\ProcedurePrintDefault{CK56070C}


\subsection{Appendizitis (Blinddarmentzündung)}
\label{topic:appendizitis}%




\TextSlide{{\TxBeschreibung}, {\TxSymptome} und Verlauf}
{%
\DefinitionInPara{Als \T{Appendizitis} bezeichnet
man die Entzündung des
\EE{Wurmfortsatzes} (\EE{Appendix} \Z{vermiformis}) des
Blinddarms. Umgangssprachlich wird oft (fälschlicherweise)
der Begriff \T{Blinddarmentzündung} verwendet.}
Im Vordergrund der akuten Appendizitis
steht anfänglich  ein meist \EE{diffuser
Schmerz} im Bereich der
\EE{Nabelgegend}, der
innerhalb weniger Stunden in den \EE{rechten Unterbauch}
wandert. Übelkeit, Brechreiz und Verstopfung können
das Krankheitsbild begleiten. Meist ist die Körpertemperatur
bis 38\DegC{} leicht erhöht. 
Unbehandelt kann es zu einer
Eiterung und einem \E{Durchbruch} in die Bauchhöhle, einer
anschließenden Bauchfellentzündung und einem akuten Abdomen
kommen (\FullPrettyRef{topic:akutes-abdomen}).
Der Verlauf einer akuten Appendizitis ist jeweils durch
individuelle Faktoren wie Lebensalter und Allgemeinzustand
des Patienten geprägt. Nicht selten sind die Symptome jedoch
völlig uncharakteristisch und evtl. sogar untypisch, sodass
eine Diagnose selbst mit erweiterten Untersuchungsmethoden
schwierig sein kann.

}{
\begin{itemize}
  \item Schmerzen im re. Unterbauch
  \item \E{Druck}schmerzen im re. Unterbauch
  \item \E{Loslass}schmerz li. Bauch
  \item Abwehrspannung
  \item Übelkeit, Erbrechen
  \item Evtl. Fieber
  \item Evtl. uncharakteristisch oder untypisch
\end{itemize}
}{}{}{}





% M %%%%%%%%%%%%%%%%%%%%%%%%%%%%%%%%%%%%%%%%%%%%%%%%%%% M %
% Akute Appendizitis
\ProcedurePrintDefault{CK35080C}



\subsection{Gallenkolik}
\label{topic:gallenkolik}%
\index{Kolik!Gallen-}%
\index{Gallenkolik}%





\TextSlide{\TxBeschreibung} 
{%
\DefinitionInPara{Eine \T{Gallenkolik} zeichnet
sich durch kolikartige Schmerzen infolge einer Blockade der
Gallenwege z.\,B. durch Gallensteine aus.} \Z{Vgl.
Kolik: \FullPrettyRef{topic:kolik}}. Gallensteine
verursachen eine Verstopfung des Gallengangs und 
folglich einen Rückstau von
Gallenflüssigkeit in die Leber. Dabei wird der Gallengang
gedehnt und es kommt zu starken Dehnungsschmerzen.} 
{
\begin{itemize}
  \item Verstopfung der Gallenwege (Stein)
  \item Rückstau und Dehnung
  \item Kolikartiger Schmerz
\end{itemize}
}{}{}{}

\SlideCompensationNoParskip{1}

\TextSlide{\TxSymptome}
{%
Durch den Rückstau und der Dehnung des Gallengangs kommt es
zu starken, auf- und abschwellenden, krampfartigen
(\EE{kolikartigen}) \EE{Schmerzen} im \EE{rechten
Oberbauch}, klassischerweise mit Ausstrahlung in die rechte
Schulterregion. Begleitendes Fieber ist möglich.

Der rechte Oberbauch ist äußerst druckschmerzhaft, der
Patient ist sehr \EE{unruhig} und agitiert.

\SlideCompensationNoParskip{2}
}{
\begin{itemize}
  \item Kolikartige Schmerzen im re. Oberbauch
  \item Nach fettreichen Mahlzeiten
  \item Evtl. Gelbfärbung des Augenweiß (Gelbsucht)
  \item Evtl. Übelkeit, Erbrechen
\end{itemize}
}{}{}{}


% M %%%%%%%%%%%%%%%%%%%%%%%%%%%%%%%%%%%%%%%%%%%%%%%%%%% M %
% Gallenkolik
\ProcedurePrintDefault{CK80811C}


% \clearpage % TEMP temp clearpage
\subsection{Gastrointestinale
Blutungen und Ösophagusvarizenblutung}
\label{topic:oesophagusvarizenblutung}%
\label{topic:gastrointestinale-blutung}%


\TextSlide{\TxBeschreibung}
{%
\DefinitionInPara{\T{Gastrointestinale Blutungen} sind
Blutungen im Magen-Darm-Trakt.}
\DefinitionInPara{\T{Ösophagusvarizenblutungen} sind
Blutungen aus erweiterten Gefäßen (\enquote*{Krampfadern}
in der Speiseröhre.} Im Verlauf des Magen-Darm-Traktes kann es an verschiedenen
Stellen zu Blutungen kommen, die auch sehr schwer ausfallen
können. Die Unterscheidung nach der Stelle der Blutung ist
wesentlich:
\begin{itemize}
  \item Blutungen im \EE{oberen Teil} des Verdauungstraktes können
  durch
  \T[Oesophagusvarizen@Ösophagusvarizen]{Ösophagusvarizen}
  verursacht werden: Bei Patienten
  mit einer fortgeschrittenen Lebererkrankung (Leberzirrhose, 
  \ref{topic:leberzirrhose}) kann es zur Bildung von 
  krampfaderartigen Gefäßerweiterungen in der Speiseröhre
  kommen. Diese Gefäße können massive Blutungen verursachen,
  die tödlich enden können. Auch Blutungen im Rahmen eines
  Magengeschwürs sind möglich, aber selten dermaßen stark
  ausgeprägt.
  \item \Z{Blutungen im \EE{mittleren Teil} des Magen-Darm-Traktes
  können durch diverse Erkrankungen und Tumore verursacht
  werden.}
  \item \Z{Blutungen im \EE{Endteil} kommen oft durch Hämorrhoiden
  oder Tumore zu Stande.}
\end{itemize}




}
{
\begin{itemize}
    \item Oberer Teil: Ösophagusvarizzen
    
    (Oft bei
    Lebererkrankungen, hoher Blutverlust möglich)
    \item \Z{Diverse Erkrankungen, Tumore}
\end{itemize}
}{}{}{}

\TextSlide{\TxSymptome}
{
Bei Blutungen im \EE{oberen Magen-Darm-Trakt}
(Speiseröhre, Magen) kommt es je nach Stärke zu
schwallartigem Erbrechen  von rotem, frischen Blut
(\Z{\I{Hämatemesis}}) oder häufiger zum Erbrechen von
\II[Erbrechen!Kaffeesatz-artiges]{Kaffeesatz-artigem}
Mageninhalt: Durch die Magensäure kommt es zur Veränderung
des Blutes, es wird bräunlich, ähnlich dem Kaffeesatz im
Kaffeefilter. Bei Blutungen im \EE{mittleren Teil} des
Verdauungstraktes färbt sich der Stuhl tief schwarz, man
spricht vom sog. \T{Teerstuhl} (\T{Melena}. Bei Butungen im
\EE{Endteil} wird rötliches Blut ausgeschieden. 

}{
\begin{itemize}
  \item Oberer Teil: Erbrechen von
  \begin{itemize}
    \item Blut im Schwall
    \item Kaffeesatz-artiger Mageninhalt
  \end{itemize} 
  \item Mittlerer Teil: Teerstuhl
  \item Endteil: Blutiger Stuhl
\end{itemize}
}{}{}{}


\Fazit{Der Patient kann aufgrund des Blutverlustes akut
vital bedroht sein, es besteht die Gefahr eines
hypovolämischen Schocks.}


% M %%%%%%%%%%%%%%%%%%%%%%%%%%%%%%%%%%%%%%%%%%%%%%%%%%% M %
% Blutungen des Verdauungstraktes
\ProcedurePrintDefault{CK92020C}



\subsection{Magen-Darm-Grippe, Gastroenteritis,
Lebensmittelvergiftung}
\label{topic:gastroenteritis}%
\label{topic:lebensmittelvergiftung}%


\TextSlide{{\TxBeschreibung} und Ursachen}
{%
\DefinitionInPara{Die \T{Magen-Darm-Grippe} ist eine
chronisch oder akut verlaufende Entzündung der Magen- und der
Dünndarmschleimhaut. \Lang{Syn.} \I{Brechdurchfall},
\Z{\I{gastrointestinaler Infekt}, \I{Gastroduodenitis}}.}
Die Gastroenteritis wird durch Erreger wie Bakterien oder
Viren hervorgerufen. Es gibt völlig harmlose aber auch
lebensbedrohliche Verlaufsformen, dies ist abhängig vom
Erreger und letztendlich auch vom Alter des Patienten. Durch
den Durchfall und das Erbrechen kann es zu einem maßgeblichen
Verlust von Wasser und Elektrolyten kommen. Wenn der Patient
nicht ausreichend trinken kann besteht die Gefahr der
Austrocknung (\I[Exsikkose!bei gostrointestinalem
Infekt]{Exsikkose}). 
}{
\begin{itemize}
  \item Vielfältig
  \item Häufig Infektionen
\end{itemize}
}{}{}{}


\Text{\TxSymptome}
{
Klassische Zeichen einer Magen-Darm-Grippe sind Erbrechen
(\Z{\I{Emesis}}) und Durchfall (\Z{\I{Diarrhoe}}). Dazu hat
der Patient meist starke, kaum lokalisierbare
Bauchschmerzen. Infolge des 
Erbrechens und des Durchfalls kann es zu einem erheblichen 
Wasser- und Elektrolytverlust und Symptomen der Exsikkose
(\ref{topic:exsikkose}) kommen. 

}{
\begin{itemize}
    \item Erbrechen 
    \item Durchfall 
    \item Bauchschmerzen
    \item Exsikkosezeichen
\end{itemize}
}{}{}{}

\TextSlide{Besonders gefährdet}
{%
Besonders betroffen davon sind Kleinkinder, da sie
normalerweise einen höheren Wasseranteil am Körpergewicht
haben.  Die Symptome sind klassisch: Oberbauchkrämpfe,
die ziemlich heftig sein können, Übelkeit, Erbrechen, Durchfälle.
Exsikkose-Gefahr! 
}{
\begin{itemize}
    \item (Klein-)Kinder
    \item Alte Menschen
\end{itemize}
}{}{}{}


% M %%%%%%%%%%%%%%%%%%%%%%%%%%%%%%%%%%%%%%%%%%%%%%%%%%% M %
% Gastroenteritis
\ProcedurePrintDefault{CA09091C}



\subsection{Flüssigkeitsmangel: Exsikkose}
\label{topic:exsikkose}%


\TextSlide{{TxBeschreibung} und Ursachen}
{%
\DefinitionInPara{Als \T{Exsikkose}
(\T{Dehydratation}, \I{Austrocknung}) wird ein hochgradiger
Mangel an Flüssigkeit im Körper bezeichnet.}
Grundsätzlich wird (wurde) entweder \E{zu wenig Flüssigkeit
zugeführt} oder \E{zu viel Flüssigkeit an die Außenwelt
abgegeben} (eine Kombination aus beidem ist natürlich auch
möglich). Dafür gibt es jeweils sehr verschiedene Ursachen:

Besonders ältere Menschen haben oft ein mangelndes
Durstgefühl und daher einen \EE{verminderten Trinkantrieb}. Auch durch die
Einnahme von \EE{Drogen} (Ecstacy) kann auf den Durst
\enquote{vergessen} werden.
Andererseits kann es zu einem erhöhtren
\EE{Flüssigkeitsverlust} durch \E{Erbrechen}, \E{Durchfall}
sowie \E{starkes Schwitzen} kommen. Bei \EE{Fieber} wird
auch vermehrt Feuchtigkeit abgeatmet, und natürlich können
auch \EE{Medikamente} entwässern.
}{
\begin{itemize} %GABS: umgebaut
  \item Mangel an Flüssigkeit
  \item Flüssigkeitszufuhr ${\downarrow}$
  \begin{compactitem}
    \item Verminderter Trinkantrieb
  \end{compactitem}
  \item Flüssigkeitsverlust ${\uparrow}$
  \begin{compactitem}
    \item Erbrechen, Durchfall
    \item Starkes Schwitzen
    \item Fieber (Feuchtigkeit wird vermehrt abgeatmet)
    \item Medikamente
  \end{compactitem}
\end{itemize}
}{}{}{}

Die genannten Ursachen sind natürlich nur ein kleiner Teil
von sehr vielen Möglichkeiten.

\TextSlide{\TxSymptome}
{%
Der Patient wirkt allgemein \E{geschwächt bis apathisch}, oft
\EE{verwirrt}, klagt über mehr oder minder heftiges
Durstgefühl, ev. über Wadenkrämpfe und Erbrechen. Die Haut
ist auffallend trocken, \EE{Falten sind abhebbar und bleiben
bestehen}. \EE{Mundhöhle und Zunge sind trocken}, evtl.
borkig oder mit eingetrocknetem Schleim belegt. Die Augen
sind charakteristischerweise tiefliegend, die Nase wirkt
auffallend spitz. Der RR ist mitunter sehr nieder, die HF
eher hoch. \E{Der Patient ist bis zur Klärung der Situation
nüchtern zu belassen!}
}{
\begin{itemize}
    \item RR ${\downarrow}$
    \item stehende Hautfalten
    \item trockene, belegte Zunge
    \item Elektrolytentgleisung
    \item \EE{Verwirrtheit, Schwindel, Apathie}
\end{itemize}
}{}{}{}



% \clearpage % TEMP temp clearpage
\section{Weniger häufige oder nicht-so-gut-erkennbare Erkrankungen}

% \minisec{Querverweise}
% 
% \begin{itemize}
%     \item Hepatitis: \prettyref{topic:hepatitis}
% \end{itemize}




\subsection{Magen- und Zwölffingerdarmgeschwür}
\label{topic:gastritis}

\TextSlide{{\TxBeschreibung} und Ursachen}
{%
\DefinitionInPara{Ein \T{Magengeschwür} bzw. ein
\T{Zwölffingerdarmgeschwür} ist eine
Entzündung der Magen- bzw. Zwölffingerdarmschleimhaut mit
Substanzdefekt. \Z{\Lang{Syn.} \I{Gastritis},
\I{Duodenitis}}.} Die häufigste Ursache für ein
Magengeschwür ist eine chronische Infektion mit einem
bestimmten Bakterium\ZFootnote{Helicobacter
pylori} \citep{Malfertheiner2004}. Bestimmte
Medikamente\footnote{Medikamente und Magengeschwür:
Betroffen sind sog. nicht-steroidale Antirheumatika (NSAR),
z.\,B.
Voltaren{\texttrademark} (Diclofenac),
Aspirin{\texttrademark} (Acetylsalicylsäure, ASS),etc.)}
schädigen die Magenschleimhaut und können bei
häufiger Einnahme zu einem Geschwür führen. Auch andere
Stoffe können auf chemisch-toxischen Weg zu einer
entsprechenden Schädigung führen. \Z{Ein
Gallensäurenrückfluß (Reflux) bzw. Autoimmunerkrankungen
können ebenso zu einem Geschwür führen.} 
}{
\begin{itemize}
    \item Bakteriell 
    \item Medikamente
    \item Chemisch-toxisch
    \item \Z{Gallensäurenrückfluß}
    \item \Z{Autoimmunerkrankungen}
\end{itemize}
}{}{}{}

\TextSlide{\TxSymptome}
{Der Patient klagt über einen
\EE{stechend-schneidenden Schmerz}, welcher sehr stark sein
kann. Daneben kommt es zu Appetitlosigkeit, Völlegefühl,
Übelkeit und Erbrechen. Wenn gleichzeitig eine \EE{Blutung}
besteht, dann sind \II[Teerstuhl!Gastritis]{Teerstühle}
(schwarz, übelriechend) oder \EE{Kaffeesatz-artiges
Erbrechen} zu beobachten. 
}{
\begin{itemize}
    \item Schmerzen (Oberbauch), Übelkeit, Erbrechen
    \item Kaffeesatz-artiges Erbrechen, Teerstuhl
\end{itemize}
}{}{}{}


% M %%%%%%%%%%%%%%%%%%%%%%%%%%%%%%%%%%%%%%%%%%%%%%%%%%% M %
% Gastritis, Duodenitis
\ProcedurePrintDefault{CK29091C}




\subsection{Pankreatitis}

\label{topic:pankreatitis}



\TextSlide{{\TxBeschreibung} und Ursachen}
{%
\index{Gallensteine!Pankreatitis}%
\index{Alkoholmißbrauch!Pankreatitis}%
\index{Bauchspeicheldrüsenentzündung}%
%%%
\DefinitionInPara{Eine \T{Pankreatitis} ist eine Entzündung
der Bauchspeicheldrüse (Pankreas). Sie kann
sowohl \E{akut}, als auch \E{chronisch} verlaufen.}
Die häufigsten Ursachen für eine Pankreatitis sind
\EE{Gallensteinleiden} und \EE{Alkoholmissbrauch}
\citep{SiewertChirurgie8-07.14-Pankreas,ScholzNotfallmedizin2}.
Daneben gibt es noch viele, allerdings nicht so häufige
Ursachen, oft bleibt die Ursache auch gänzlich unbekannt
\citep{HeroldInnereMedizin2005}.

\Speech{Warum können Gallensteine eine Pankreatitis
verursachen?} Bei vielen Menschen mündet der
Pankreasausführungsgang, welcher Verdauungsenzyme in de
Zwölffingerdarm abgibt, nicht direkt in den Dünndarm,
sondern in den unteren Abschnitt des Gallengangs.
Verschließt nun ein Stein den Gallengang unterhalb der
Einmündung, so staut sich Galle und Pankreassaft in die
Bauchspeicheldrüse zurück. Es kann sogar dazu kommen, dass
sich das Pankreas durch die Verdauungsenzyme selber
an-/verdaut.

}
{\begin{itemize}
  \item Gallensteinleiden
  \item Alkoholmißbrauch
  \item \ldots
\end{itemize}}
{}
{}
{}

\TextSlide{\TxSymptome}
{Die Symptome sind oft unspezifisch, eine
Erstdiagnosestellung ist präklinisch normalerweise nicht
möglich. Charakteristisch ist bestenfalls ein
\II{gürtelförmiger
Oberbauchschmerz}\index{Oberbauchschmerz,gürtelförmiger}.
Die häufigsten Begleitsymptome sind Übelkeit und Erbrechen,
sowie Fieber. \citep{HeroldInnereMedizin2005}

\bigskip
}
{\begin{itemize}
  \item Oberbauchschmerzen, oft gürtelförmig
  \item Übelkeit, Erbrechen
  \item Fieber
\end{itemize}}
{}
{}
{}

\TextSlide{\TxKomplikationen}
{Je nach Schwere und Dauer der Entzündung kann es zum
Selbstverdauen der Bauchspeicheldrüse bzw. zu einer
\E{Ruptur} kommen. In Folge tritt ein
\I[Akutes Abdomen!bei Pankreatitis]{akutes Abdomen} auf
(\FullPrettyRef{topic:akutes-abdomen}) und der Patient kann
kann einen \I[Schock!septischer!bei Pankreatitis]{septischen
Schock} erleiden (\FullPrettyRef{topic:schock-septischer}).

}
{\begin{itemize}
  \item Ruptur
  \item Akutes Abdomen
  \item Septischer Schock
\end{itemize}}
{}
{}
{}


% M %%%%%%%%%%%%%%%%%%%%%%%%%%%%%%%%%%%%%%%%%%%%%%%%%%% M %
% Akute Pankreatitis
\ProcedurePrintDefault{CK859X0C}


\subsection{Mesenterialinfarkt}
\label{topic:mesenterialinfarkt}

\Text
{\TxBeschreibung}
{\DefinitionInPara{Ein \T{Mesenterialinfarkt} ist eine
Ischämie des Darms durch Verschluss einer Darmarterie.
\Z{\Lang{Syn.} \I{Darminfarkt}}.} Die den
Darm versorgenden Arterien verlaufen im Mesenterium\ZFootnote{Das
\I[Mesenterium!Mesenterialinfarkt]{Mesenterium}
(\I{Darmgekröse}) ist die Aufhängung von Teilen des
Dünndarms (\FullPrettyRef{topic:mesenterium}).
Darin verlaufen venöse und arterielle Gefäße zum und vom Darm
(\I{Mesenterialgefäße} mit \I[Mesenterialarterie]{Mesenterialarterien}).}, 
daher leitet sich der Name der Erkrankung ab.


Die Folge eines Mesenterialinfarktes ist das \EE{Absterben
des entsprechenden Darmabschnitts}, dadurch kommt es zu
einer Perforation mit einer Bauchfellentzündung, Sepsis und
einem Versterben des Patienten.
Er ist somit ein \I[Ischämie!Mesenterialinfarkt]{ischämischer}
Notfall. 
Mesenterialinfarkte treten gehäuft bei Patienten mit
erhöhtem thromboembolischen Risikofaktoren auf, z.\,B. bei
Patienten mit
\I[Vorhofflimmern!Mesenterialinfarkt]{Vorhofflimmern}
(\FullPrettyRef{topic:vorhofflimmern}).

}
{}
{}
{}
{}


\TextSlide{\TxSymptome}
{Die Symptome sind eher unspezifisch. Der Patient klagt oft
(aber nicht immer) über \EE{abdominelle Schmerzen}, welche
oft dumpf und diffus sind. Es kann für einige Stunden ein
\E{stilles Intervall} folgen, in welchem die Schmerzen eher
gering sind. 

Nach ungefähr \VonBis{12}{24}\Hour* wird ein
\EE{Darmverschluss} bemerkbar und die Darmwand verliert
zunehmend an Dichtigkeit, es kann zu einem \EE{Durchbruch}
(Ruptur) kommen. Es folgt eine \EE{Bauchfellentzündung} und
ein \II[Akutes Abdomen!beim Mesenterialinfarkt]{akutes
Abdomen} (\FullPrettyRef{topic:akutes-abdomen}). Es kann zu
\EE{blutigen, übel riechenden Stuhlabgängen} kommen.
\cite{SiewertChirurgie8-06-Gefaesschirurgie}
% 
% \bigskip

}
{\begin{itemize}
  \item Abdominalschmerzen
  \item Übelkeit, Erbrechen
  \item $\sim$ 6\Hour* stilles Intervall
  \item Nach $\sim$ \VonBis{12}{24}\Hour* Durchbruch und
  Bauchfellentzündung, Ileus
  \item Übelriechende, blutige Stuhlabgänge
\end{itemize}}
{}
{}
{}


% M %%%%%%%%%%%%%%%%%%%%%%%%%%%%%%%%%%%%%%%%%%%%%%%%%%% M %
% Mesenterialinfarkt
\ProcedurePrintDefault{CK55001C}



% \clearpage % TEMP temp clearpage
\subsection{Bauchfellentzündung}
\label{topic:bauchfellentzuendung}
\label{topic:peritonitis}


\TextSlide{\TxBeschreibung}{%
\DefinitionInPara{Eine \T{Bauchfellentzündung}
(\I{Peritonitis}) ist eine Entzündung bzw. Keimbesiedelung
des Bauchfells.}
Durch die Entzündung bzw. Infektion  des Bauchfelles
kommt es zu Symptomen eines Akuten Abdomens. Oft ist die Perforation
eines Hohlorganes  (Magendurchbruch, Darmperforation,
Blinddarmdurchbruch, \ldots) Ursache der Entzündung. 
Eine Bauchfellentzündung führt üblicherweise zu einer
schweren  körperweiten Entzündungsreaktion
(\I[Sepsis!Bauchfellentzünfdung]{Sepsis},
\FullPrettyRef{topic:sepsis}).

}{
\begin{itemize}
  \item Entzündung bzw. Infektion des Bauchfells
\end{itemize}
}{}{}{}

\TextSlide{Ursache}{
\begin{itemize}
    \item Perforation eines Hohlorganes (Darm, Magen,
    Gallenblase, \ldots)
    \item Fremdkörper
\end{itemize}
}{
\begin{itemize}
    \item Perforation eines Hohlorganes
    \item Infektion über Blut- / Lymphwege oder eingewandert
    über Fremdkörper
\end{itemize}
}{}{}{}

\TextSlide{\TxSymptome}{
\begin{itemize}
    \item Symptome eines Akuten Abdomens
    \item Schmerzen (diffus)
    \item Evtl. reflektorischer Darmverschluss
    \item Evtl. Sepsis- und/oder Schockzeichen
\end{itemize}

}{
\begin{itemize}
    \item Symptome eines Akuten Abdomens
    \item Schmerzen (diffus)
    \item Evtl. reflektorischer Darmverschluss
    \item Evtl. Sepsis- und/oder Schockzeichen
\end{itemize}
}{}{}{}


% M %%%%%%%%%%%%%%%%%%%%%%%%%%%%%%%%%%%%%%%%%%%%%%%%%%% M %
% Bauchfellentzündung
\ProcedurePrintDefault{CK65090C}


% \Z{
% \subsection{Reserviert \A{Sonstige -- nicht akute --
% Erkrankungen}}
% \subsubsection{Reserviert \A{Chronisch entzündliche
% Darmerkrankungen}} }
% 
\ChapterEnd
